\subsection{Rigidity at the original irreducible points}
\label{sec:support-type-one-rigidity}

Thorne's adjoint Selmer theorem applies over the auxiliary CM field
already constructed.  This avoids descending the automorphic
representation before proving rigidity.  The local conditions and the
action of complex conjugation must still be retained exactly.

Write $K/\mathbf Q$ for the original imaginary quadratic field, $\eta$
for its character, and $J$ for the $3$ by $3$ antidiagonal matrix with
antidiagonal entries one.  Let $\Psi=\varepsilon^e\chi$ be the original
even character, with $e\ge1$ and $\chi$ of finite order; it is not the
bare cyclotomic character.  The residual character is $\psi$.
The field $K$ is unramified at the odd coefficient prime $\ell$.
The original local decomposition group $D$ has an element
\[
       x\in D\cap G_K,\qquad \psi(x)^2\ne1.             \tag{TR1}
\]
The mark $x$ need not be in inertia.  This distinction retains, in
particular, the original conductor-one construction at $\ell=3$.
Write $I$ for the inertia subgroup of $D\cap G_K$.

\begin{lemma}[Irreducibility after restriction]
\label{lem:support-type-one-strong-irreducibility}
Let $E/\mathbf Q_\ell$ be finite, and let
$\rho:G_K\to\mathrm{GL}_3(E)$ be continuous and absolutely irreducible.
Suppose it has an invariant integral lattice with residual
semisimplification $\psi\oplus1\oplus\psi^{-1}$.  Suppose
$\rho|_{D\cap G_K}$ preserves a full flag whose ordered inertia
characters are $\Psi,1,\Psi^{-1}$, and assume (TR1).
Then $\rho$ remains absolutely irreducible on every open subgroup of
$G_K$.  If $L/K$ is finite and $L(\mu_{\ell^\infty})$ is its
cyclotomic extension, it is also absolutely irreducible on
$G_{L(\mu_{\ell^\infty})}$.
\end{lemma}

\begin{proof}
Work over $\overline{\mathbf Q}_\ell$.  It suffices first to consider
an open normal subgroup $N$ of $G_K$, by replacing any open subgroup
with its normal core.  If $\rho|_N$ is reducible, Clifford theory in
dimension three gives either a scalar action of $N$, or a direct sum
of three lines permuted transitively by $G_K$.  For completeness,
restriction is semisimple: the radical of the finite-dimensional image
algebra of $N$ is stable under conjugation by $G_K$, so its image on
the irreducible space is zero, by nilpotence.  The isotypic summands
are permuted transitively and have equal dimensions and multiplicities.
The alternatives follow since three is prime.

There is $u\in I\cap N$ for which $\Psi(u)$ has infinite order.
Indeed the cyclotomic image of inertia over the finite local field is
open in $\mathbf Z_\ell^\times$; imposing membership in $N$ and
killing the finite character $\chi$ still leaves an open subgroup.
The three eigenvalues $\Psi(u),1,\Psi(u)^{-1}$ are distinct.
This rules out the scalar alternative.  In the other alternative
the three $N$-lines are exactly the three eigenspaces of $\rho(u)$.
The local ordinary flag is an ordered flag of these lines.  Every
element of $D\cap G_K$ permutes the lines and preserves this flag,
so fixes each line.

A transitive subgroup of $S_3$ contains a $3$-cycle.  Choose $g\in G_K$
acting by such a cycle.  Since $x$ fixes the three lines,
$\rho(gx^m)$ has trace zero for every $m\ge0$.  Reduction on the
invariant lattice gives
\[
  \alpha\beta^m+1+(\alpha\beta^m)^{-1}=0,
  \qquad \alpha=\psi(g),\quad\beta=\psi(x).            \tag{TR2}
\]
For $m=0,1,2$ the three values $\alpha\beta^m$ are distinct, since
$\beta^2\ne1$.  Equation (TR2) makes all three roots of
$T^2+T+1$, a contradiction in every characteristic, including three.
This proves the assertion for open subgroups.

Here are also the details for the infinite cyclotomic extension.
Let $\mathcal G$ be the Zariski closure of $\rho(G_L)$.
The inverse image of its identity component $\mathcal G^0$ is a
closed subgroup of finite index in the compact group $G_L$, hence is
open.  Therefore $\mathcal G^0$ acts irreducibly.  Its unipotent
radical has a nonzero fixed vector and is normal; irreducibility makes
it act trivially on the whole space.  As $\mathcal G^0$ is a matrix
group, that radical is trivial.  Thus $\mathcal G^0$ is reductive.
Its connected center acts by scalars, and
$\mathcal G^0=Z(\mathcal G^0)^0[\mathcal G^0,\mathcal G^0]$.
Consequently its derived subgroup acts irreducibly as well.
The Zariski closure of $\rho([G_L,G_L])$ contains this derived
subgroup: apply the commutator morphism to the dense image of
$G_L\times G_L$.  The cyclotomic quotient is abelian, so
$[G_L,G_L]\subseteq G_{L(\mu_{\ell^\infty})}$.  This proves the
last assertion, with no assertion about equality of local and global
cyclotomic intersections.
\end{proof}

\begin{theorem}[The original type-I rigidity]
\label{thm:support-type-one-rigidity}
For either original admitted datum, retain any allowed local ring
$U$ and its associated normalized global deformation ring $S$.
Let $f:S\to\overline{\mathbf Q}_\ell$ be a local characteristic-zero
point, so $|f(\mathfrak m_S)|<1$, and suppose its representation is
type I: its matrix part $\rho$ on $G_K$ is absolutely irreducible.
Then
\[
                 H_f^1(\mathbf Q,\operatorname{ad}r_f)=0.         \tag{TR3}
\]
Here $\operatorname{ad}r_f$ is the original polarized adjoint on
$M_3(\overline{\mathbf Q}_\ell)$: $G_K$ acts by conjugation by
$\rho$, and the normalized complex conjugation $c$ acts by
$X\mapsto-JX^{\mathsf t}J$.

Every original algebra-valued dual-number lift of $f$ has tangent
cocycle $X-g\cdot X$ for some matrix $X$.  Thus the entire original
rigidity requirement at type-I points holds.  No residual
irreducibility, residual image-size condition, or uniform bound on an
integral Selmer group is imposed or inferred.
\end{theorem}

\begin{proof}
\emph{The actual point and its automorphy.}
The finite-neighbourhood argument of
Proposition~\ref{prop:selmer-tangent-comparison} puts the point over
the integers $O_E$ of a finite coefficient extension.  The local
ordered-flag lemma~\ref{lem:support-integral-ordinary-flag} applies
to its actual local ring, and the residual representation has the
specified diagonal characters.  Lemma~\ref{lem:support-type-one-strong-irreducibility}
therefore applies.

Use the actual auxiliary $L/L^+$ and its original quotient and frame,
with the full admission, $\mathrm{Ok}$, $\mathrm{Big}$ and seed
conditions already proved.  The normalized-lift transfer
Theorem~\ref{thm:support-finite-end} sends $r_f$ to the rank-three
normalized representation over this same $O_E$.  Restriction to
$G_L$ is irreducible by the lemma.  The original finite-point
Theorem~\ref{thm:support-finite-points}
therefore gives a point of the actual ordinary Hecke quotient whose
determinant is that of
\[
                  \rho_L'=\rho|_{G_L}\,\varepsilon^{-1}.        \tag{TR4}
\]
Its ordered inertia characters are
$\Psi^{1-i}\varepsilon^{-1}$, $i=0,1,2$.  In the original
Hodge--Tate convention their labels are $1-e(1-i)$, strictly
increasing because $e\ge1$.  Their remaining factors have finite
order, so this is an arithmetic point in the precise sense of the
ordinary-family theorem.  The full arithmetic-point classicality in
Theorem~\ref{thm:selmer-ordinary-families}, applied to this
irreducible representation, makes (TR4) the representation of a
regular algebraic conjugate-self-dual cuspidal representation
$\Pi$ of $\mathrm{GL}_3(\mathbb A_L)$.  The normalized-lift and
Hecke determinant identities identify the actual representation,
not merely its residual determinant.  Enlarging coefficients to
contain the original coefficient datum changes neither $L$ nor its
level.  No automorphy descent from $L$ to $K$ is used.

\emph{The adjoint representation.}
The prescribed polarization is
\[
 \rho(cgc^{-1})=J\rho(g)^{-\mathsf t}J.
\]
After (TR4) its multiplier on $G_L$ is $\varepsilon^{-2}$;
on $G_{L^+}$ the polarized multiplier is
$\eta_{L/L^+}\varepsilon^{-2}$.  In the usual polarized group
$\mathcal G_3$, the chosen complex conjugation is $(J,1)j$,
where $j$ acts on the Lie algebra by negative transpose.
Consequently its adjoint action is exactly
\[
                         X\longmapsto-JX^{\mathsf t}J.           \tag{TR5}
\]
Scalar Tate twisting does not alter the adjoint action.  Furthermore,
any polarization of the irreducible $\rho_L'$ is a scalar multiple
of this one, by Schur's lemma; that scalar does not alter (TR5).
Thus the adjoint in Thorne's theorem is precisely
$(\operatorname{ad}r_f)|_{G_{L^+}}$, including its one-dimensional
trace summand, on which complex conjugation acts by minus one.

Lemma~\ref{lem:support-type-one-strong-irreducibility} also makes
$(\rho_L')|_{G_{L(\mu_{\ell^\infty})}}$ irreducible.  Thorne's
Theorem~6.1 now gives
\[
             H_f^1(L^+,\operatorname{ad}\rho_L')=0
                                      \cite{support-type-one-thorne}.
\]
Restriction preserves both the unramified local condition away from
$\ell$ and the Bloch--Kato finite condition at $\ell$.  It is
injective on global $H^1$ in characteristic zero: corestriction
composed with restriction is multiplication by $[L^+:\mathbf Q]$.
This degree may be divisible by $\ell$, but it is invertible in the
coefficient field.  Equation (TR3) follows.

\emph{The actual dual-number condition.}
For an original dual-number lift, the same finite-neighbourhood
argument gives a continuous tangent cocycle over one finite
coefficient extension.  At $\ell$ use the established
crystalline-component comparison,
Proposition~\ref{prop:selmer-crystalline-components}, on the
original selected saturated components.  Its nonresonance argument
uses the original factors of $\delta$, including the exponent-one
case.  It proves crystallinity of the entire dual-number lift after
the finite local extension killing the original finite inertial
character.  Thus its self-extension class is in $H_f^1$ there.
Restriction injectivity for
$B_{\mathrm{cris}}\otimes\operatorname{ad}r_f$ descends that finite
condition to the original local group.  This uses the full adjoint
coefficient module, not an inclusion into a larger representation.
In particular ordinary triangularity alone is not being used as a
crystalline criterion.

Away from $\ell$ the original equations kill the open inertia
subgroup on which $\eta$ and the finite character $\chi$ are trivial;
here $\Psi=\chi$ on inertia.
The tangent class therefore factors on inertia through a finite
group.  Averaging over that group makes it a coboundary, as required
by the unramified local condition.  The tangent class consequently
belongs to the group in (TR3) and is zero.  Its primitive over the
finite coefficient field gives the asserted original coboundary
after extension to $\overline{\mathbf Q}_\ell$.
\end{proof}

The theorem closes the type-I part of the original two-sided rigidity
obligation.  Type-II points have additional rank-two representation
summands, which are not adjoint representations and are not covered by
Thorne's theorem.  Their vanishing, and the moving-family integral
Selmer estimate, remain separate requirements.
